\chapter{Fonetica e pronuncia del Pāli}

Il Pāli è la lingua scritturale originale del Buddhismo Theravāda. Era una lingua
parlata, strettamente imparentata con il Sanscrito, senza un proprio sistema di scrittura.
Con l'emergere di forme scritte, queste hanno utilizzato gli alfabeti di altre lingue
(ad es. Devanagari, Singalese, Birmano, Khmer, Tailandese, Romano). La trascrizione in caratteri romani
qui utilizzata si pronuncia come in italiano, con le seguenti precisazioni:

\section{Vocali}

\begin{tabular}{@{} ll @{}}

\begin{tabular}{@{} ll @{}}
  Corte & Lunghe\\
  \textbf{a} as in \prul{a}bout &
  \textbf{ā} as in f\prul{a}ther\\
  \textbf{i} as in h\prul{i}t &
  \textbf{ī} as in mach\prul{i}ne\\
  \textbf{u} as in p\prul{u}t &
  \textbf{ū} as in r\prul{u}le\\
  & \textbf{e} as in gr\prul{e}y\\
  & \textbf{o} as in m\prul{o}re\\
\end{tabular} &

\parbox[b][][t]{0.5\linewidth}{%
Eccezioni: \textbf{e} e \textbf{o} diventano suoni brevi nelle sillabe
che terminano in consonante. Si pronunciano quindi come in 't\prul{e}tto' e
'c\prul{o}tto', rispettivamente.} \\

\end{tabular}

\section{Consonanti}

\textbf{c} come in \prul{c}ento (come \prul{ch} ma non aspirata)

\textbf{ṁ, ṅ} come \prul{ng} in ma\prul{ng}o

\textbf{ñ} come \prul{ni} in co\prul{ni}o

\textbf{v} più morbida che la nostra \prul{v}; più vicina a \prul{w}

\textbf{cc} è una doppia \prul{c} come in Fibona\prul{cc}i, mai pronunciata come in co\prul{cc}o.

\subsection{Consonanti aspirate}

\textbf{bh ch dh ḍh gh jh kh ph th ṭh}

Queste notazioni a due lettere con \prul{h} indicano un suono aspirato e arioso,
distinto dal suono duro e secco della singola consonante. Dovrebbero essere
considerate come un'unica lettera.


Tuttavia, le altre combinazioni con \textbf{h}, cioè \textbf{lh, mh, ñh,} e
\textbf{vh}, contano come due consonanti (ad esempio nelle parole Pāli
'ji\textbf{vh}ā' o 'mu\textbf{ḷh}o').

\ifaivedition
\clearpage
\fi

\subsection{Esempi}

\textbf{th} come \prul{t} in \prul{t}avolo.

\textbf{ph} as \prul{p} in \prul{p}iatto. (Mai pronunciata come in `\prul{ph}oto'.)

Questi suoni sono distinti dal suono duro e secco della singola consonante, ad es.
\textbf{th} come in '\prul{Th}omas' (non come nell'inglese '\prul{th}in') o \textbf{ph} come
in '\prul{p}uff' (non come in '\prul{f}oto').

\subsection{Consonanti retroflesse}

\textbf{ḍ ḍh ḷ ṇ ṭ ṭh}

Queste consonanti retroflesse non hanno equivalenti in italiano. Si pronunciano
arricciando la punta della lingua all'indietro contro il palato.

\section{Tecnica di canto}

Una volta compreso il sistema di pronuncia Pāli e la seguente
tecnica di canto, è possibile cantare un testo in Pāli a prima vista
con il ritmo corretto.

\textbf{Le sillabe non accentate} terminano con una \textbf{a, i} o
\textbf{u} breve. Tutte le altre sillabe sono accentate. Le sillabe accentate durano
il doppio del tempo di quelle non accentate --- un po' come due battiti in una battuta
musicale rispetto a uno. Questo è ciò che conferisce al canto il suo particolare
ritmo.

\begin{centering}

{\setlength{\tabcolsep}{1.8pt}%
\begin{tabular}{ccc c ccccc c ccccc c ccccccc}
BUD & · & DHO & \hsp & SU & · & SUD & · & DHO & \hsp & KA & · & RU & · & ṆĀ & \hsp & MA & · & HAṆ & · & ṆA & · & VO\\
 1  &   & 1   &      & ½  &   & 1   &   & 1   &      & ½  &   & ½  &   & 1  &      & ½  &   & 1   &   & ½  &   & 1\\
\end{tabular}%
}

\end{centering}

Due dettagli importanti nella separazione delle sillabe:

\textbf{1.} Sillabe con lettere doppie vengono divise in questo modo:
\begin{centering}

\begin{minipage}{0.8\linewidth}
\begin{multicols}{2}
\setlength{\tabcolsep}{1.8pt}%

\begin{tabular}{rrcccl}
     & A & · & NIC & · & CA   \\
     & ½ &   &  1  &   & ½    \\
(not & A & · & NI  & · & CCA) \\
     & ½ &   & ½   &   & ½    \\
\end{tabular}

\columnbreak

\begin{tabular}{rrcccl}
     & PUG & · & GA  & · & LĀ \\
     &  1  &   &  ½  &   &  1 \\
(not & PU  & · & GGA & · & LĀ)\\
     &  ½  &   &  ½  &   &  1 \\
\end{tabular}

\end{multicols}
\end{minipage}

\end{centering}

Vengono sempre enunciate separatamente, ad es. \textbf{dd} in 'uddeso' come
in 'a\textbf{dd}io', o \textbf{gg} in 'maggo' come in 'a\textbf{gg}ancio'.

\textbf{2.} \textbf{Consonanti aspirate} come \textbf{bh, dh} etc.
contano come una singola consonante e non vengono divise (Quindi
\textbf{am·hā·kaṁ}, ma \textbf{sa·dham·maṁ}, non \textbf{sad·ham·maṁ};
un altro esempio: \textbf{Bud·dho} e non \textbf{Bu·ddho}).

Una pronuncia precisa e una corretta separazione delle sillabe sono
particolarmente importanti per chi è interessato a imparare il Pāli e a
comprendere e memorizzare il significato dei Sutta e di altri canti,
altrimenti il significato verrà distorto.

\textbf{Un esempio per illustrare questo punto:}


La parola Pāli '\textbf{sukka}' significa 'luminoso'; '\textbf{sukkha}' significa
'secco'; '\textbf{sukha}' --- 'felicità'; '\textbf{suka}' --- 'pappagallo' e
'\textbf{sūka}' --- 'Lo stelo di una spiga d'orzo'.

Quindi, se si canta '\textbf{sukha}' con una '\textbf{k}' invece di una
'\textbf{kh}', si canterebbe 'pappagallo' invece di 'felicità'.

Una regola generale per comprendere la pratica del canto è
ascoltare attentamente ciò che il leader e il gruppo cantano e seguire,
mantenendo la stessa intonazione, tempo e velocità. Tutte le voci dovrebbero fondersi
in una sola.

\section{Punteggiatura, segni tonali e pause in questa edizione}

[Le parentesi quadre] indicano le parti solitamente cantate solo dal leader, ma
le usanze di canto differiscono nei vari monasteri.

La barra / indica le variazioni delle forme maschili o femminili a seconda
della persona che canta, o le forme singolari e plurali quando si canta
da soli o in gruppo.

I segni di cantillazione indicano cambiamenti di intonazione, di solito un tono intero sopra o sotto:

\begin{tabular}{llll}
High tone: & no꜓ble & Long low tone: & ho꜖mage\\
Low tone: & ble꜕ssed & Long mid tone: & \prul{guides}\\
\end{tabular}

\section{Nota sulla sillabazione nel testo}

Come aiuto alla comprensione, alcune delle parole Pāli più lunghe nel testo sono state
suddivise con un trattino nelle parole da cui sono composte. Questo non influisce
in alcun modo sulla pronuncia.

% In order to not suggest unintended pauses in the flow of the chanting, we have
% omitted all punctuation marks (commas, periods, colon and semicolon), although
% for rendering the meaning of the phrases accurately, they would be required.
% The line breaks indicate that a short breathing pause is inserted.

% End of pali-phonetics-and-pronunciation.tex
